\documentclass[10pt,a4paper]{amsart}
\usepackage[margin=19mm]{geometry}
\usepackage{amsmath,amssymb,mathtools}
\usepackage[scr=rsfs,cal=euler]{mathalfa}
\usepackage{xfrac}
\usepackage[unicode,hidelinks,bookmarks=false]{hyperref}
\newcommand{\ie}{i.e.\ }
\newcommand{\cf}{cf.\ }
\newcommand{\resp}{respectively\ }
\newcommand{\Subsection}{\subsection}
\providecommand{\nobreakdash}{\nobreak\hbox{-}}
\setlength{\emergencystretch}{2em}
\multlinegap0pt
\usepackage{bm}
\usepackage{bbm}
\usepackage{dsfont}
\usepackage{xspace}
\usepackage{cancel}
\usepackage{mathtools}
\usepackage[shortlabels]{enumitem}
\usepackage{amsthm}
\makeatletter
\@ifundefined{thm}{\newtheorem{thm}{Thm}}{}
\@ifundefined{lemma}{\newtheorem{lemma}[thm]{Lemma}}{}
\makeatother
\def \r{{\mathbb{R}}}
\renewcommand{\Gamma}{\varGamma}
\title[The Dirichlet problem for a family of totally degenerate dif]{The Dirichlet problem for a family of totally degenerate differential operators}
\author{Maria Manfredini, Mirco Piccinini, Sergio Polidoro}
\date{Source: arXiv:2106.12048v3 (CC BY 4.0)}
\begin{document}
\maketitle

\noindent\emph{Source and licence.} The Dirichlet problem for a family of totally degenerate differential operators. Version arXiv:2106.12048v3 (CC BY 4.0), distributed under the Creative Commons Attribution 4.0 International licence, as recorded in the arXiv licence metadata for this exact version and shown on the abstract page https://arxiv.org/abs/2106.12048v3. Full rights notice: licenses/arxiv-2106.12048-CC-BY-4.0.txt.\par\smallskip

\noindent\textit{Editorial scope.} Counted unit: the Lemma whose source label is \texttt{rapporto} in the article (source0.tex lines 570--576, source label \texttt{rapporto}), delivered together with the article's own complete original proof of it (source0.tex lines 577--667, one \texttt{proof} environment of 91 lines). No mathematics is written, repaired or re-derived: statement and proof are the article's own text line for line. Required context retained uncounted: def:normC (lines 339--341); confrhomnonhom (lines 359--362); fund-sol (lines 503--510); lemma:dilation-fundamental (lines 513--525); eq:commutator (lines 529--531); matinv (lines 533--536); eq:notation (lines 562--568). Every definition, display and label of those lines is retained unchanged. Omitted from this package and not redistributed: 1--338, 342--358, 363--502, 511--512, 526--528, 532--532, 537--561, 569--569, 668--1745. Nothing inside a retained excerpt is omitted or altered: every retained body is delivered exactly as the article's lines are written, so no omission receipt of any kind accompanies this package. Citations: the retained excerpts carry no citation command at all, so no bibliography entry of the article is transcribed and no reference is redistributed. Reference adaptation: none was needed and none was made; every reference inside the delivered unit resolves inside it, so no unresolved reference or citation is produced. Printed slips kept literal: no printed word, symbol, hypothesis or number of the counted statement or of its proof was altered. Layout and class: the article's own class is \texttt{amsart} with packages including lmodern, textcomp, latexsym, amssymb, amsbsy, amstext, esint, enumerate, cancel, cite, bbm, mathrsfs, ulem, lipsum; the wrapper is the lane's pinned \texttt{amsart} editorial wrapper at 10pt on A4, which restates the theorem-like environments the retained text needs and adds the article's own macro definitions that the retained text needs (\texttt{\textbackslash Gamma}, \texttt{\textbackslash r}), with extra wrapper packages (bm, bbm, dsfont, xspace, cancel, mathtools, enumitem). The delivered numbering is the unit's own. Assets: no retained span contains a figure environment, a graphics-inclusion command, a PDF-page inclusion, a table or a TikZ picture; no image file is copied into this package, the package declares no asset dependency and no image file is redistributed with it. Licence: version arXiv:2106.12048v3 of this article is distributed under the Creative Commons Attribution 4.0 International licence, as recorded in the arXiv licence metadata for this exact version and shown on the abstract page https://arxiv.org/abs/2106.12048v3. A full rights notice is delivered at licenses/arxiv-2106.12048-CC-BY-4.0.txt. Characters: every retained excerpt is pure ASCII, so the delivered engine, XeLaTeX with its default Latin Modern text font, reports no missing character.
\begin{equation}\label{def:normC}
|x|_{C}^2 := \frac{1}{4}\langle C^{-1}(-1,0) x,x\rangle.
\end{equation}

\begin{eqnarray}\label{confrhomnonhom}
&& \sigma\min{\left\{\left|x\right|^\frac{1}{1+2\vartheta}, \left|x\right|^{\frac{1}{2(\kappa+\vartheta)+1}}\right\}}\notag\\
&& \qquad \quad \leq\left\|x\right\|\leq (\kappa+1)\max{\left\{\left|x\right|^\frac{1}{1+2\vartheta}, \left|x\right|^{\frac{1}{2(\kappa+\vartheta)+1}}\right\}}\qquad \forall\,x\in\r^N,
\end{eqnarray}

\begin{equation} \label{fund-sol}
	\Gamma(z;\zeta):= 
    \begin{cases}
        \frac{(4\pi)^{-N/2}}{\sqrt{\textup{det}C(\tau,t)}}
	e^{-\frac{1}{4}\langle C^{-1}(\tau,t)( x- e^{-(t-\tau)B}\xi ) , x- e^{-(t-\tau)B}\xi \rangle} & \quad \text{if} \ t > \tau\\
    0 & \quad \text{if}\ t \le \tau.
        \end{cases}
\end{equation}

\begin{lemma}\label{lemma:dilation-fundamental}
The following properties of the fundamental solution $\Gamma$ in \eqref{fund-sol} hold true:
\begin{enumerate}[(i)]
 \item   For any $(x,t),(\xi,\tau) \in \mathbb{R}^{N+1}$ 
         \[
         \Gamma(x,t;\xi,\tau) = \Gamma(x-e^{-(t-\tau)B}y,t;\xi-y,\tau) \quad \forall y \in \r^N.
         \]
 \item   For any $(x,t),(\xi,\tau) \in \mathbb{R}^{N+1}$ and any $r>0$ it holds
    \[
   \Gamma(x,t;\xi,\tau) =  r^{Q}\Gamma(D_r x, r^2 t; D_r \xi, r^2 \tau).
    \]
\end{enumerate}
\end{lemma}

\begin{equation}\label{eq:commutator}
e^{-r^2sB}D_r=D_r e^{-sB}\qquad \forall r>0,\, \forall s\in \r.
\end{equation}

\begin{lemma}\label{matinv}
For $0>t>\tau$ we have the following matrix inequality
$$e^{-tB^T}C^{-1}(\tau,t)e^{-tB}\geq C^{-1}(\tau,0).$$
\end{lemma}

\begin{equation}\label{eq:notation}
\begin{split}
\mu & :=\frac{-t}{-\tau}\in(0,1),\\
 {M}(z_{\rm o},z) & :=\left|D_{\frac{1}{\sqrt{-t}}}(x-e^{-tB}x_{\rm o})\right|,\\
 \text{and}\quad {M}(z_{\rm o},\zeta)& :=\left|D_{\frac{1}{\sqrt{-\tau}}}(\xi-e^{-\tau B}x_{\rm o})\right|
 \end{split}
\end{equation}

\begin{lemma}\label{rapporto}
Fix $z_{\rm o}=(x_{\rm o},0) \in \r^{N+1}$. There exists a positive constant $c$ such that, for any $z=(x,t),\zeta=(\xi,\tau)$ with $0>t>\tau$ and $\mu\leq\min{\{\frac{1}{2},\frac{\sigma^2}{(\kappa+1)^2}\}}$, we have
\[
\frac{\Gamma(z,\zeta)}{\Gamma(z_{\rm o},\zeta)}\leq\left(\frac{1}{1-\mu}\right)^{\frac{Q}{2}}e^{c\sqrt{\mu}M(z_{\rm o},z)M(z_{\rm o},\zeta)}.
\]
where $\mu$ and $M(\cdot)$ are both defined in \eqref{eq:notation}.
\end{lemma}

\begin{proof} 
Applying the  transformation in Lemma \ref{lemma:dilation-fundamental} we obtain that 
\begin{eqnarray*}
  \Gamma(z,\zeta)  & = & (t-\tau)^{-\frac{Q}{2}} \Gamma(D_{\frac{1}{\sqrt{t-\tau}}}(x-e^{-(t-\tau)B}\xi),\frac{t}{t-\tau};0,\frac{\tau}{t-\tau})\\
  & = & \frac{c(t-\tau)^{-\frac{Q}{2}}}{\sqrt{\det C(\frac{\tau}{t-\tau},\frac{t}{t-\tau})}}e^{-\frac{1}{4}\langle C^{-1}(\tau,t)(x-e^{-(t-\tau)B}\xi),x-e^{-(t-\tau)B}\xi\rangle}\\
   & = & \frac{c(t-\tau)^{-\frac{Q}{2}}}{\sqrt{\det C(-\frac{1}{1-\mu},-\frac{\mu}{1-\mu})}}e^{-\frac{1}{4}\langle C^{-1}(\tau,t)(x-e^{-(t-\tau)B}\xi),x-e^{-(t-\tau)B}\xi\rangle}\,,
\end{eqnarray*}
and, similarly, we have
\begin{eqnarray*}
    \Gamma(z_{\rm o},\zeta) &=&\frac{c(-\tau)^{-\frac{Q}{2}}}{\sqrt{\det C(-1,0)}}e^{-\left|D_{\frac{1}{\sqrt{-\tau}}}(x_{\rm o}-e^{\tau B}\xi)\right|_C^2}\,,
\end{eqnarray*}
recalling the definition of $|\cdot|_C$ in \eqref{def:normC}. 

Then, since 
\[
\frac{(t-\tau)^{-\frac{Q}{2}}\sqrt{\det C(-1,0)}}{(-\tau)^{-\frac{Q}{2}}\sqrt{\det C(-\frac{1}{1-\mu},-\frac{\mu}{1-\mu})}} = \left(\frac{1}{1-\mu}\right)^{\frac{Q}{2}}\,,
\] 
the only thing we need to control is the exponential term. For this, start noticing that
\begin{eqnarray}\label{eq:conto}
    x-e^{-(t-\tau)B}\xi & = & x - e^{-(t-\tau)B}(\xi -e^{-\tau B}x_{\rm o}+ e^{-\tau B}x_{\rm o})\notag\\
    & = & x - e^{-(t-\tau)B}e^{-\tau B}x_{\rm o} - e^{-(t-\tau)B}(\xi -e^{-\tau B}x_{\rm o})\notag\\
    & = &  x - e^{-tB}x_{\rm o}- e^{-(t-\tau)B}(\xi -e^{-\tau B}x_{\rm o})\,,
\end{eqnarray}
where we have used that $(e^{-tB})^{-1}=e^{t B}$. Then, by  \eqref{eq:conto} we obtain
\begin{eqnarray}\label{exponent}
&&-\frac{1}{4}\langle C^{-1}(\tau,t)(x-e^{-(t-\tau)B}\xi),x-e^{-(t-\tau)B}\xi\rangle\notag\\
&& \quad = -\frac{1}{4}\langle C^{-1}(\tau,t)(x-e^{-tB}x_{\rm o}),x-e^{-tB}x_{\rm o}\rangle\notag\\
&& \qquad -\frac{1}{4}\langle C^{-1}(\tau,t)e^{-(t-\tau)B}(\xi -e^{-\tau B}x_{\rm o}),e^{-(t-\tau)B}(\xi -e^{-\tau B}x_{\rm o})\rangle\nonumber\\
&& \qquad + \frac{1}{2}\left(\langle C^{-1}(\tau,t)(x-e^{-tB}x_{\rm o}),x-e^{-tB}x_{\rm o}\rangle\right)^\frac{1}{2}\notag\\
&& \qquad \times \left( \langle C^{-1}(\tau,t)e^{-(t-\tau)B}(\xi -e^{-\tau B}x_{\rm o}),e^{-(t-\tau)B}(\xi -e^{-\tau B}x_{\rm o})\rangle\right)^\frac{1}{2}\notag\\ 
&& \quad \leq -\frac{1}{4}\langle C^{-1}(\tau,t)e^{-(t-\tau)B}(\xi -e^{-\tau B}x_{\rm o}),e^{-(t-\tau)B}(\xi -e^{-\tau B}x_{\rm o})\rangle\\
&& \qquad + \frac{1}{2}\left(\langle C^{-1}(\tau,t)(x-e^{-tB}x_{\rm o}),x-e^{-tB}x_{\rm o}\rangle\right)^\frac{1}{2}\notag\\
&& \qquad \times \left( \langle C^{-1}(\tau,t)e^{-(t-\tau)B}(\xi -e^{-\tau B}x_{\rm o}),e^{-(t-\tau)B}(\xi -e^{-\tau B}x_{\rm o})\rangle\right)^\frac{1}{2} \,, \notag
\end{eqnarray}
since $C^{-1}(\cdot,\cdot)$ is a positive definite matrix.  Moreover, Lemma \ref{matinv} yields that for any $y \in \r^N$
 $$
 \left\langle C^{-1}(\tau,0)e^{\tau B}y,e^{\tau B}y\right\rangle-\left\langle C^{-1}(\tau,t) e^{-(t-\tau)B}y,e^{-(t-\tau)B}y\right\rangle\leq 0.
 $$
 By using this and \eqref{exponent} we get
\begin{eqnarray}\label{eq:exponent2}
&&\left|D_{\frac{1}{\sqrt{-\tau}}}\left(x_{\rm o}-e^{\tau B}\xi\right)\right|_C^2-\frac{1}{4}\langle C^{-1}(\tau,t)(x-e^{-(t-\tau)B}\xi),x-e^{-(t-\tau)B}\xi\rangle\notag\\
 && \quad \leq  \frac{1}{4}\langle C^{-1}(\tau,0)e^{\tau B}\left(\xi - e^{-\tau B}x_{\rm o}\right),e^{\tau B}\left(\xi - e^{-\tau B}x_{\rm o}\right)\rangle \notag\\
 && \qquad -\frac{1}{4}\langle C^{-1}(\tau,t)e^{-(t-\tau)B}(\xi -e^{-\tau B}x_{\rm o}),e^{-(t-\tau)B}(\xi -e^{-\tau B}x_{\rm o})\notag\\
&& \qquad + \frac{1}{2}\left(\langle C^{-1}(\tau,t)(x-e^{-tB}x_{\rm o}),x-e^{-tB}x_{\rm o}\rangle\right)^\frac{1}{2}\notag\\
&& \qquad \times \left( \langle C^{-1}(\tau,t)e^{-(t-\tau)B}(\xi -e^{-\tau B}x_{\rm o}),e^{-(t-\tau)B}(\xi -e^{-\tau B}x_{\rm o})\rangle\right)^\frac{1}{2} \notag\\
&& \quad \leq \frac{1}{2}\left(\langle C^{-1}(\tau,t)(x-e^{-tB}x_{\rm o}),x-e^{-tB}x_{\rm o}\rangle\right)^\frac{1}{2}\\
&& \qquad \times \left( \langle C^{-1}(\tau,t)e^{-(t-\tau)B}(\xi -e^{-\tau B}x_{\rm o}),e^{-(t-\tau)B}(\xi -e^{-\tau B}x_{\rm o})\rangle\right)^\frac{1}{2}\notag
\end{eqnarray}
We are going to bound all the above terms in \eqref{eq:exponent2} separately. We first have
\begin{eqnarray*}
&& \langle C^{-1}(\tau,t)(x-e^{-tB}x_{\rm o}),x-e^{-tB}x_{\rm o}\rangle\\*[0.5ex]
&& \quad =\left\langle C^{-1}\left(-\frac{1}{\mu},-1\right)D_{\frac{1}{\sqrt{-t}}}(x-e^{-tB}x_{\rm o}),D_{\frac{1}{\sqrt{-t}}}(x-e^{-tB}x_{\rm o})\right\rangle.
\end{eqnarray*}
Now, let us denote with $\left\|A\right\|$ the operator norm of a matrix $A$ (i.e. its biggest eigenvalue for symmetric matrices). By \eqref{confrhomnonhom}, for any vector $v$ with $\left|v\right|=1$ we get
\begin{eqnarray*}
\min{\left\{\left|D_{\sqrt{\mu}}v\right|^\frac{1}{1+2\vartheta},\left|D_{\sqrt{\mu}}v\right|^{\frac{1}{2(\kappa+\vartheta)+1}}\right\}}\leq\frac{1}{\sigma}\sqrt{\mu}\left\|v\right\| & \leq & \frac{\kappa+1}{\sigma}\sqrt{\mu}\max{\left\{\left|v\right|^\frac{1}{2\vartheta+1},\left|v\right|^{\frac{1}{2(\kappa+\vartheta)+1}}\right\}}\\
& = & \frac{\kappa+1}{\sigma}\sqrt{\mu}.
\end{eqnarray*}

From $\mu\leq \frac{\sigma^2}{(\kappa+1)^2}$ we then deduce $\left|D_{\sqrt{\mu}}v\right|\leq \left(\frac{\kappa+1}{\sigma}\right)^{2\vartheta +1}\mu^{\frac{1}{2}(2\vartheta +1)} \leq \left(\frac{\kappa+1}{\sigma}\right)^{2\vartheta +1}\sqrt{\mu} $ since $\mu \in (0,1)$. Hence, since by definition $\mu$ is also less than $\frac{1}{2}$,  for any $|v|=1$ we obtain that
\begin{eqnarray*}
\left\langle C^{-1}\left(-\frac{1}{\mu},-1\right)v,v \right\rangle&=&\left\langle C^{-1}(-1,-\mu)D_{\sqrt{\mu}}v,D_{\sqrt{\mu}}v\right\rangle\\
& \leq &  \left\|C^{-1}(-1,-\mu)\right\|\left|D_{\sqrt{\mu}}v\right|^2\\
&\leq& \left(\frac{\kappa+1}{\sigma}\right)^{4\vartheta +2}\left\|C^{-1}\left(-1,-\mu\right)\right\|\mu\\
& \leq &  \left(\frac{\kappa+1}{\sigma}\right)^{4\vartheta +2}\left\|C^{-1}\left(-1,-\frac{1}{2}\right)\right\|\mu\,,
\end{eqnarray*}
which gives recalling \eqref{eq:notation}
\begin{eqnarray*}
&& \left\langle C^{-1}(\tau,t)(x-e^{-tB}x_{\rm o}),x-e^{-tB}x_{\rm o}\right\rangle\notag\\
&& \qquad\quad \leq \left(\frac{\kappa+1}{\sigma}\right)^{4\vartheta +2}\left\|C^{-1}\left(-1,-\frac{1}{2}\right)\right\|\mu M(z_{\rm o},z)^2.
\end{eqnarray*}
On the other hand, by the commutation property \eqref{eq:commutator}, we get
\begin{eqnarray*}
&& \left\langle C^{-1}(\tau,t)e^{-(t-\tau)B}(\xi -e^{-\tau B}x_{\rm o}),e^{-(t-\tau)B}(\xi -e^{-\tau B}x_{\rm o})\right\rangle\notag\\
&& \quad =  \left\langle C^{-1}(-1,-\mu)D_{\frac{1}{\sqrt{-\tau}}}e^{-(t-\tau)B}(\xi -e^{-\tau B}x_{\rm o}),D_{\frac{1}{\sqrt{-\tau}}}e^{-(t-\tau)B}(\xi -e^{-\tau B}x_{\rm o})\right\rangle\\
&& \quad \leq\left\|C^{-1}(-1,-\mu)\right\|\left|D_{\frac{1}{\sqrt{-\tau}}}e^{-(t-\tau)B}(\xi -e^{-\tau B}x_{\rm o})\right|^2\notag\\
&& \quad = \left\|C^{-1}(-1,-\mu)\right\|\left|e^{-(1-\mu)B}D_{\frac{1}{\sqrt{-\tau}}}(\xi -e^{-\tau B}x_{\rm o})\right|^2\\
&& \quad \leq \left\|C^{-1}(-1,-\mu)\right\|\left\|e^{-(1-\mu)(B+B^T)}\right\|\left|D_{\frac{1}{\sqrt{-\tau}}}(\xi -e^{-\tau B}x_{\rm o})\right|^2\\
&& \quad \leq \left\|C^{-1}\left(-1,-\frac{1}{2}\right)\right\|\left\|e^{-(1-\mu)(B+B^T)}\right\|\left|D_{\frac{1}{\sqrt{-\tau}}}(\xi -e^{-\tau B}x_{\rm o})\right|^2.
\end{eqnarray*}
Since $0<\mu\leq\frac{1}{2}$, the term $\left\|e^{-(1-\mu)(B+B^T)}\right\|$ is bounded from above by a universal constant $c_0^2$. Thus we have
\begin{eqnarray*}
&& \left\langle C^{-1}(\tau,t)e^{-(t-\tau)B}(\xi -e^{-\tau B}x_{\rm o}),e^{-(t-\tau)B}(\xi -e^{-\tau B}x_{\rm o})\right\rangle\notag\\
&& \qquad\quad \leq c_0^2 \left\|C^{-1}\left(-1,-\frac{1}{2}\right)\right\|M(z_{\rm o},\xi)^2.
\end{eqnarray*}
Therefore
\[
\frac{\Gamma(z,\zeta)}{\Gamma(z_{\rm o},\zeta)}\leq\left(\frac{1}{1-\mu}\right)^{\frac{Q}{2}}e^{c\sqrt{\mu} M(z_{\rm o},z) M(z_{\rm o},\xi)}\,,
\]
which gives the thesis.
\end{proof}

\end{document}
